% Kandydat restrukturyzacji; nie jest aktywnym rozdziałem Atlasu.
\section{Lokalność, nośniki i przedłużanie dystrybucji}
\providecommand{\atlasref}[2]{#2}

\subsection{Nośnik i lokalność}

\begin{definicja}[Zerowanie się dystrybucji na zbiorze otwartym]
\label{def:distribution-zero19}
Niech $T\in\mathcal D'(\Omega)$, a $U\subseteq\Omega$
będzie otwarty. Dla $U=\varnothing$ przyjmujemy
$\mathcal D(U;\K)=\{0\}$ i $T|_U=0$.
Dla niepustego $U$ i $\varphi\in\mathcal D(U;\K)$ oznaczmy przez
$e_U\varphi$ jej przedłużenie zerem na $\Omega$.
Jest ono gładkie, gdyż nośnik $\varphi$ jest zwarty
i zawarty w $U$. Na każdym etapie $\mathcal D_K(U)$
odwzorowanie $e_U$ zachowuje seminormy $p_m$ i ma obraz
w $\mathcal D_K(\Omega)$, więc jest ciągłe.
Z Twierdzenia~\ref{thm:LF-universal18} wynika ciągłość
$e_U:\mathcal D(U;\K)\to\mathcal D(\Omega;\K)$.
Definiujemy ograniczenie $T|_U$ przez
$\langle T|_U,\varphi\rangle=\langle T,e_U\varphi\rangle$.
Jest ono dystrybucją jako złożenie ciągłych odwzorowań liniowych.
Piszemy $T|_U=0$, jeśli powyższa wartość jest zerowa
dla każdej $\varphi\in\mathcal D(U;\K)$.
Nośnikiem $T$ jest domknięty względnie zbiór
\[\supp T=\Omega\setminus\bigcup\{U\subseteq\Omega:
U\text{ jest otwarty i }T|_U=0\}.\]
\end{definicja}

\begin{lemat}[Lokalność dystrybucji]
\label{lem:distribution-local19}
Niech $T\in\mathcal D'(\Omega)$ i
$\varphi\in\mathcal D(\Omega)$ spełnia
$\supp\varphi\cap\supp T=\varnothing$. Wtedy
$\langle T,\varphi\rangle=0$.
\end{lemat}
\begin{proof}
Jeśli $\varphi=0$, teza jest oczywista.
Niech $L=\supp\varphi\ne\varnothing$.
Każdy $x\in L$ ma otoczenie otwarte $U_x\subseteq\Omega$,
na którym $T$ się zeruje. Wybierzmy kulę $B_x$
o środku $x$ i dodatnim promieniu, z $\overline B_x\subseteq U_x$.
Ze zwartości $L$ wybieramy skończone pokrycie
$B_{x_1},\ldots,B_{x_N}$.
Z Lematu~\ref{lem:cutoffs17} wybieramy nieujemne
$\eta_j\in C_c^\infty(U_{x_j};\R)$ równe $1$
w otoczeniu $\overline B_{x_j}$ i przedłużamy je zerem.
Niech $h=\sum_{j=1}^N\eta_j$ i $O=\{x\in\Omega:h(x)>0\}$.
Zbiór $O$ jest otwarty i zawiera $L$.
Na $O$ definiujemy $\varphi_j=\eta_j\varphi/h$,
a na $\Omega\setminus O$ przyjmujemy $\varphi_j=0$.
Przedłużenie jest gładkie, ponieważ każdy punkt spoza $O$
ma otoczenie rozłączne z domkniętym zbiorem $L$,
na którym $\varphi_j$ jest identycznie zerowa.
Mamy $\supp\varphi_j\subseteq L\cap\supp\eta_j\Subset U_{x_j}$
oraz $\sum_j\varphi_j=\varphi$ na $O$ i poza $O$.
Zatem $\langle T,\varphi_j\rangle=0$ dla każdego $j$.
Sumując, otrzymujemy tezę.

W szczególności $T$ zeruje się na otwartym zbiorze
$\Omega\setminus\supp T$, gdyż każda funkcja testowa
na tym zbiorze, przedłużona zerem, ma nośnik rozłączny
z $\supp T$. Jest to więc największy otwarty zbiór
zerowania się $T$.
\end{proof}

\begin{stwierdzenie}[Nośniki przykładów]
\label{prop:support-examples19}
Dla $a\in\Omega$ i $\alpha\in\N_0^d$ mamy
$\supp\delta_a=\supp T_{a,\alpha}=\{a\}$.
Dla $f\in L^1_{\mathrm{loc}}(\Omega)$ nośnik $T_f$
jest dopełnieniem największego otwartego podzbioru
$U\subseteq\Omega$, na którym $f=0$ prawie wszędzie.
\end{stwierdzenie}
\begin{proof}
Gdy otwarty $U\subseteq\Omega$ nie zawiera $a$,
przedłużenie każdej funkcji z $\mathcal D(U)$ znika
w otoczeniu $a$, zatem wszystkie jej pochodne w $a$ są zerowe.
Jeśli $a\in U$, wybieramy $\eta\in C_c^\infty(U;\R)$
równą $1$ w otoczeniu $a$. Funkcja
$\varphi(x)=(x-a)^\alpha\eta(x)$ spełnia
$\partial^\alpha\varphi(a)=\alpha!\ne0$.
Stąd żadne otoczenie $a$ nie jest zbiorem zerowania
$T_{a,\alpha}$, także dla $\alpha=0$.
Dla $T_f$ mamy $(T_f)|_U=T_{f|_U}$ z definicji ograniczenia.
Jeśli $f=0$ prawie wszędzie w $U$, to $(T_f)|_U=0$.
Odwrotnie, ze Stwierdzenia~\ref{prop:regular-injective19}
na $U$ wynika, że $(T_f)|_U=0$ pociąga $f=0$ prawie wszędzie w $U$.
Największy taki otwarty zbiór istnieje na mocy Lematu
o lokalności, a jego dopełnienie jest nośnikiem $T_f$.
\end{proof}

\begin{twierdzenie}[Zwarty nośnik implikuje skończony rząd]
\label{thm:compact-support-order19}
Jeżeli $T\in\mathcal D'(\Omega;\K)$ i
$S=\supp T\Subset\Omega$, to istnieją
$m\in\N_0$, $C\ge0$ i $K\Subset\Omega$ takie, że
\[|\langle T,\varphi\rangle|\le C p_{K,m}(\varphi)
\quad(\varphi\in\mathcal D(\Omega)),\qquad
p_{K,m}(\varphi)=\max_{|\alpha|\le m}
\sup_{x\in K}|\partial^\alpha\varphi(x)|.\]
W szczególności $T$ ma skończony rząd na całej $\Omega$.
Ponadto $T$ przedłuża się w sposób jednoznaczny do
ciągłego funkcjonału liniowego na $C^\infty(\Omega;\K)$.
\end{twierdzenie}
\begin{proof}
Jeśli $S=\varnothing$, Lemat~\ref{lem:distribution-local19}
daje $T=0$. Możemy wtedy wybrać $m=0$, $C=0$
i dowolny niepusty zwarty $K\Subset\Omega$.
Przedłużeniem jest funkcjonał zerowy; jego jednoznaczność
wynika z gęstości $\mathcal D(\Omega)$ w $C^\infty(\Omega)$
na mocy Wniosku~\ref{cor:density-cinf17}.
W przeciwnym przypadku z Lematu~\ref{lem:cutoffs17}
wybieramy $\chi\in C_c^\infty(\Omega;\R)$ równą $1$
na pewnym otoczeniu $S$. Wtedy
$\supp((1-\chi)\varphi)\cap S=\varnothing$, więc z
Lematu~\ref{lem:distribution-local19} wynika
$\langle T,\varphi\rangle=\langle T,\chi\varphi\rangle$.
Niech $K=\supp\chi$. Z Twierdzenia~\ref{thm:distribution-estimate19}
istnieją $m,C_0$ takie, że
$|\langle T,\psi\rangle|\le C_0p_m(\psi)$ dla
$\psi\in\mathcal D_K$. Reguła Leibniza daje
$p_m(\chi\varphi)\le 2^m p_m(\chi)p_{K,m}(\varphi)$,
ponieważ wszystkie pochodne $\chi$ znikają poza $K$.
Otrzymujemy żądane oszacowanie z $C=C_0 2^m p_m(\chi)$.
Dla każdego zwartego $L\Subset\Omega$ i $\varphi\in\mathcal D_L$
mamy $p_{K,m}(\varphi)\le p_m(\varphi)$,
więc jeden rząd $m$ wystarcza dla wszystkich $L$.
Dla $u\in C^\infty(\Omega)$ definiujemy
$\widetilde T(u)=\langle T,\chi u\rangle$.
Odwzorowanie to jest liniowe, a powyższe oszacowanie
iloczynu obowiązuje również dla $u\in C^\infty(\Omega)$.
Zatem $|\widetilde T(u)|\le C p_{K,m}(u)$,
co zapewnia ciągłość w topologii $C^\infty(\Omega)$, a wzór z poprzedniego akapitu
pokazuje zgodność z $T$ na funkcjach testowych.
Przestrzeń $\mathcal D(\Omega)$ jest gęsta w
$C^\infty(\Omega)$ w tej topologii (Wniosek~\ref{cor:density-cinf17}),
więc przedłużenie jest jednoznaczne.
\end{proof}

\subsection{Lokalność i problemy przedłużania dystrybucji}

\begin{twierdzenie}[Ograniczenie i brak automatycznego przedłużenia]
\label{thm:restriction25}
Jeżeli $U\subseteq\Omega\subseteq\R^d$ są niepustymi
zbiorami otwartymi, to odwzorowanie ograniczenia
$\mathcal D'(\Omega)\to\mathcal D'(U)$ jest liniowe
i ciągłe w topologiach słabych-* oraz silnych.
Nie dla każdej pary $U\subseteq\Omega$ odwzorowanie
ograniczenia jest surjektywne. Nawet dla przedziałów
$U=(0,1)$ i $\Omega=(-1,1)$ istnieją dystrybucje
regularne na $U$, które nie mają przedłużenia na $\Omega$.
\end{twierdzenie}
\begin{proof}
Funkcję testową $\varphi\in\mathcal D(U)$ przedłużamy
przez zero do $\Omega$ i oznaczamy przedłużenie przez
$j\varphi=\widetilde\varphi$. Jeżeli
$\supp\varphi\subseteq L\Subset U$, to $L$ jest również
zwarty w $\Omega$, a każdy punkt $\Omega\setminus U$
ma otoczenie rozłączne z $L$. Przedłużenie jest więc
gładkie także na brzegu $U$ w $\Omega$ i
$\supp(j\varphi)=\supp\varphi$, gdzie oba nośniki
utożsamiamy jako podzbiory $\Omega$.
Na każdym niezerowym etapie $\mathcal D_L(U)$
odwzorowanie $j$ ma wartości w $\mathcal D_L(\Omega)$
i zachowuje seminormy $p_{L,m}$, ponieważ pochodne
na $U$ nie zmieniają się. Na etapie zerowym ciągłość
jest oczywista. Własność uniwersalna granicy induktywnej
(Twierdzenie~\ref{thm:LF-universal18}) zapewnia ciągłość
liniowego odwzorowania $j:\mathcal D(U)\to\mathcal D(\Omega)$.

Ograniczenie $R_U T=T|_U$ definiujemy przez
$\langle R_U T,\varphi\rangle=\langle T,j\varphi\rangle$.
Jest to ciągły funkcjonał na $\mathcal D(U)$ jako
złożenie ciągłych odwzorowań liniowych.
Dla każdego ograniczonego $B\subseteq\mathcal D(U)$
mamy dokładną równość
\[
 \sup_{\varphi\in B}|\langle R_U T,\varphi\rangle|
 =\sup_{\eta\in j(B)}|\langle T,\eta\rangle|.
\]
Ponieważ $j(B)$ jest ograniczony, otrzymujemy ciągłość
silną, a dla $B$ jednoelementowego także słabą-*.
Nie zakładamy tu gęstości obrazu $j$.

Dla otwartego $V\subseteq U$ kolejne przedłużenia przez
zero dają to samo odwzorowanie co bezpośrednie
przedłużenie z $V$ do $\Omega$. Zatem
$(T|_U)|_V=T|_V$. Z definicji nośnika i
Lematu~\ref{lem:distribution-local19} wynika
$\supp(T|_U)=U\cap\supp T$, z nośnikiem po lewej
stronie rozumianym względnie w $U$: punkt $x\in U$
nie należy do tego nośnika wtedy i tylko wtedy,
gdy istnieje otwarte otoczenie $V\subseteq U$ punktu
$x$, na którym $T|_V=0$. W szczególności
$\ker R_U=\{T\in\mathcal D'(\Omega):
\supp T\subseteq\Omega\setminus U\}$.
Jeżeli przedłużenie istnieje, jego jednoznaczność
nie jest na ogół zapewniona: do dowolnego przedłużenia
można dodać element tego jądra.

Aby wykazać drugie stwierdzenie, weźmy
$U=(0,1)$, $\Omega=(-1,1)$ i dystrybucję regularną
$T_f$ wyznaczoną przez $f(x)=\exp(1/x^2)$ na $U$.
Funkcja $f$ jest gładka na $U$, a więc lokalnie
całkowalna. Ustalmy $K=[-1/2,1/2]\Subset\Omega$.
Gdyby istniało przedłużenie
$\widetilde T\in\mathcal D'(\Omega)$ z
$\widetilde T|_U=T_f$, to z
Twierdzenia~\ref{thm:distribution-estimate19}
dla tego ustalonego $K$ istniałyby $m\in\N_0$
i $C>0$ takie, że
$|\langle\widetilde T,\varphi\rangle|\le C p_{K,m}(\varphi)$
dla wszystkich $\varphi\in\mathcal D_K(\Omega)$.
Niech $\psi\in\mathcal D((-1,1))$ będzie nieujemna
i niezerowa. Przedłużamy ją przez zero na $\R$.
Dla całkowitych $n\ge3$ przyjmujemy
$r_n=1/(10n^2)$ i
$\varphi_n(x)=\psi((x-1/n)/r_n)$.
Nośnik zawiera się w przedziale
$[1/n-r_n,1/n+r_n]\subset(0,1/2)$, zatem
$\varphi_n$ jest testem zarówno na $U$, jak i na
ustalonym etapie $\mathcal D_K(\Omega)$.
Dla każdego $0\le k\le m$ mamy
$\varphi_n^{(k)}(x)=r_n^{-k}\psi^{(k)}((x-1/n)/r_n)$.
Ponieważ $r_n\le1$, otrzymujemy
\[
 p_{K,m}(\varphi_n)
 \le 10^m n^{2m}\max_{0\le k\le m}
 \sup_{\R}|\psi^{(k)}|.
\]
Z drugiej strony dla $x=1/n+r_nt$, $|t|\le1$,
mamy $0<x\le(11/10)n^{-1}$, a więc
$1/x^2\ge(100/121)n^2\ge n^2/2$.
Podstawienie w całce i nieujemność $\psi$ dają
\[
 \begin{aligned}
 \langle T_f,\varphi_n\rangle
 &=r_n\int_{-1}^1 e^{1/(1/n+r_nt)^2}\psi(t)\,
 \mathrm d\lambda_1(t)\\
 &\ge \frac{e^{n^2/2}}{10n^2}
 \int_{-1}^1\psi(t)\,\mathrm d\lambda_1(t).
 \end{aligned}
\]
Stała $c=\frac1{10}\int_{-1}^1\psi(t)\,\mathrm d\lambda_1(t)$
jest ściśle dodatnia. Ponieważ
$\langle\widetilde T,\varphi_n\rangle=
\langle T_f,\varphi_n\rangle$, z oszacowania na
$\mathcal D_K(\Omega)$ wynikałoby
$c e^{n^2/2}\le C' n^{2m+2}$ dla wszystkich $n\ge3$,
gdzie $C'$ nie zależy od $n$. Jednak
$e^{n^2/2}\ge (n^2/2)^{m+2}/(m+2)!$,
więc po podzieleniu przez $n^{2m+2}$ lewa strona
rośnie co najmniej jak dodatnia stała razy $n^2$.
Otrzymujemy sprzeczność. Brak przedłużenia dotyczy
wszystkich dystrybucji, nie tylko dystrybucji regularnych.
\end{proof}
